% SCRATCH DRAFT — Vol 2 courtroom coda, Ch20 "Frames, Clocks, and the Interval" (task vol2-coda, author: Podo).
% THIS IS NOT THE BOOK. Do not \input. Floor commit not yet banked; 20.tex untouched.
% PART V OPENER — ⭐ the THIRD HAND finally WORKED (takes the stand). Fresh-append (no old beat; body ends line 212).
%
% ⚠ FENCES:
%  - ±1 VALUE + spinor OUT. Sagnac = the THIRD proved loop residue: MAY announce PROVED (a proved failure to
%    close) but WITHHOLD the value (like Ch5/Ch16) — no number, no sign.
%  - ⭐ CONSENSUS RECALIBRATION: the third hand is NOW PRESENT, so consensus MAY rise ABOVE the Ch6 ceiling —
%    all three on the bench, the meeting ASSEMBLABLE, the correction that reconciles a bending clock being WORKED.
%    BUT the actual THREE-WAY AGREEMENT / the fine consistent under law stays RESERVED for Ch30: do NOT reach the
%    three agreeing on one number, do NOT name/reach the fine or the verdict. Advance the meeting to its WORK, not its END.
%  - Charge re-sounds (the count now carries a GEOMETRY — measured along a path by the interval / proper time).
%    Continuum foundational floor reserved. "verdict"/2nd-person/penalty-"fine"/free-will ABSENT.
%
% CONTINUITY: state carries from Ch19. Ch19 closed: "the third hand comes to the bench… the meeting draws near…
% its third reader finally on its way." Ch20 IS the third hand taking the stand.
%
% PROPOSED SEAM in books/experimentation/latex/chapters/20.tex:
%   KEEP the entire locked body through line 212 ("...the potentials that bend the record's own geometry.").
%   APPEND: \bigskip / \begin{center}\rule{0.4\textwidth}{0.4pt}\end{center} / \bigskip / then the coda below. (Fresh append.)
%
% >>>>>>>>>>>>>>>>>>>>>> APPEND BEGINS (fresh \rule + coda after the body) >>>>>>>>>>>>>>>>>>>>>>

\bigskip
\begin{center}
\rule{0.4\textwidth}{0.4pt}
\end{center}
\bigskip

\noindent The third witness takes the stand at last. Through nineteen chapters it was named and awaited --- the
reader that neither rides the wheel nor stands at the roadside, but looks down from far overhead and reads the
car by a clock of its own. Now it is here, and the first thing it says unsettles the court: there is no
privileged clock, and no privileged place from which to keep one. Every frame is a labeling, one among many, and
none of them the world's to prefer --- each reading the case has gathered was taken from somewhere, and no
somewhere is the true one.

This is the bending measure the last chapter promised, and it bends further than the court expected. Set two
clocks moving relative to one another and they will not keep the same time; ask two observers in motion which of
two distant events came first and they need not agree; what is one witness's single \emph{now} is another's
slanting slice of before-and-after. The overhead reader's clock runs at its own rate, neither the wheel's nor
the field's, so the three readings the case has gathered no longer even share a measure in which to be compared.
The frame is ours to choose; that none of them is privileged, and that the measures diverge, is the world's.
This is no figure of speech. Muons made high in the air, living at rest only a sliver of time, reach the ground
in numbers only because their own clocks, racing near the speed of light, run slow enough to stretch that sliver
across the whole descent; and atomic clocks flown around the world aboard airliners came home reading a little
behind the clocks left at rest, having ticked, in flight, a measurably smaller count. The bending is a reading
on a caesium clock, not a manner of speaking.

And yet the divergence is not chaos, for beneath the disagreeing clocks lies one thing every frame keeps the
same. Give any clock a path through events and it counts, along that path, a number all observers agree on
however they slice their time and their space: the interval, the proper time the clock itself tallies as it
goes. The count the book has carried since the first mark has, here, acquired a geometry --- and the geometry is
nothing but the measure of the count along its own path. Which frame reads it, and how the reading is cut into
time and space, is ours; the interval the clock counts is the world's, one and the same for all. Nor is any
direction privileged over another: light sent down two crossed arms and back returns in step however the whole
is turned, and that stubborn agreement is the measure declining to prefer a way to point. (Even a turning
leaves its mark: send a signal both ways around a spinning ring and the two do not come home in step --- a
proved failure to close, the third time the instrument carries a tally around a real loop and proves it comes
back charged, though what that residue is worth the book, as ever, keeps back.)

And the \emph{charge} takes on that geometry too. The charge was a count; now the count is measured along a path,
and no single standing clock reads it --- what a thing accumulates depends on the road it took to get there, and
two roads between the same two points need not tally the same. So the charge the case will collect is no longer
a number read off one dial; it is an interval counted along a path, the proper time of the thing that carried
it. The loop the charge counts is itself a path through events, and the interval is its true length. The twin who
stays home on a straight road ages more than the one who loops away and returns, because the straight path
counts the longer interval --- the most proper time between two meetings falls to the plainest route between
them. The trace
is that path-length held invariant --- the one measure that survives every witness's slicing, kept whole where
each frame's own clock would disagree. Whatever comes to be owed is measured along the road, not read from a
clock standing still.

And now, for the first time, all three witnesses are in the room together. The wheel's reading from within, the
field's from without, and the overhead reader's from above --- the three the case was built upon, present at
last, the meeting no longer a promise but a thing that can be assembled. It will not assemble itself. The third
witness reads by a clock that bends, and its measure does not fall into line with the other two of its own
accord; to bring the three onto one number, the case must first learn the correction that reconciles a bending
clock with standing ones. That reconciliation is the work the coming chapters take up --- the calibration that
lets three differently-clocked readings be made to meet. The three are on the bench; the number they must be
brought to is not yet in hand. The meeting has begun its work; it has not reached its end.

The overhead reader is only half-calibrated, though. What these chapters opened is the geometry of motion ---
clocks that slow with speed, frames with no privilege among them --- but the reader hangs high above the road,
and height, it will turn out, bends the measure as surely as speed does: a clock lifted up does not run at the
rate of one below. The next chapter turns to gravity proper --- the redshift, the delay, the pull that bends the
record's own geometry --- the other half of what the third witness needs before it can be believed. The case is
still open, its three witnesses finally assembled and not yet agreed; the court does not sit until the record is
complete.

% <<<<<<<<<<<<<<<<<<<<<< APPEND ENDS <<<<<<<<<<<<<<<<<<<<<<
